\documentclass[12pt,a4paper]{article}
\usepackage[a4paper,margin=2cm]{geometry}
\usepackage[T2A]{fontenc}
\usepackage[utf8]{inputenc}
\usepackage[russian]{babel}
\usepackage{amsmath,amssymb,mathtools}
\usepackage{array,booktabs,longtable}
\usepackage{xcolor}
\usepackage{hyperref}
\usepackage{tikz}

\hypersetup{
  colorlinks=true,
  linkcolor=blue!55!black,
  urlcolor=blue!55!black
}

\setlength{\parindent}{0pt}
\setlength{\parskip}{0.6em}
\renewcommand{\arraystretch}{1.25}

\newcommand{\errlabel}{\textcolor{red}{\textbf{Ошибка:}}}
\newcommand{\keyidea}{\textcolor{blue!60!black}{\textbf{Ключевая идея:}}}
\newcommand{\resultlabel}{\textcolor{teal!60!black}{\textbf{Итог:}}}

\begin{document}

\begin{titlepage}
\centering
\vspace*{0.22\textheight}
{\Huge\bfseries Ревью домашней работы\par}
\vspace{2.2cm}
{\Large Автор: Лёвин Артём Александрович\par}
\vspace{1cm}
{\large 23 сентября 2026 года\par}
\vfill
\begin{tikzpicture}
\draw[line width=0.7pt,blue!45!black]
  (0,0) .. controls (2,0.8) and (4,-0.8) .. (6,0)
        .. controls (8,0.8) and (10,-0.8) .. (12,0);
\end{tikzpicture}
\end{titlepage}

\newpage
\section*{Сводная таблица}
\small
\setlength{\tabcolsep}{3pt}
\begin{longtable}{>{\raggedright\arraybackslash}p{1.0cm}
                  >{\raggedright\arraybackslash}p{2.2cm}
                  >{\raggedright\arraybackslash}p{4.1cm}
                  >{\raggedright\arraybackslash}p{3.4cm}
                  >{\raggedright\arraybackslash}p{4.0cm}}
\toprule
\textbf{№} & \textbf{Тема} & \textbf{Суть ошибки} & \textbf{Ваш ответ} & \textbf{Правильный ответ} \\
\midrule
\endfirsthead
\toprule
\textbf{№} & \textbf{Тема} & \textbf{Суть ошибки} & \textbf{Ваш ответ} & \textbf{Правильный ответ} \\
\midrule
\endhead
\hyperref[task:2]{2} &
Сокращение дроби и график функции &
После верного упрощения значение $x=1$ использовано как обычная точка графика. Выколотая точка не оформлена, а график не доведён до требуемой параболы. &
Вершина $(0,-4)$, пересечений с осью $Ox$ нет, пересечение с осью $Oy$ --- $(0,-4)$; выколотая точка не указана. &
$y=-x^2-4$, $x\ne1$; вершина $(0,-4)$; пересечений с $Ox$ нет; пересечение с $Oy$ --- $(0,-4)$; выколотая точка $(1,-5)$. \\
\midrule
\hyperref[task:3]{3} &
Число пересечений с горизонтальной прямой &
Не рассмотрен промежуток $-\frac{25}{4}<m<-4$, где горизонтальная прямая имеет две общие точки с графиком. &
Одна точка при $m=-\frac{25}{4}$ и $m=-4$; две точки при $m>-4$; нет точек при $m<-\frac{25}{4}$. &
Нет точек при $m<-\frac{25}{4}$; одна точка при $m=-\frac{25}{4}$ или $m=-4$; две точки при $m>-\frac{25}{4}$, $m\ne-4$. \\
\midrule
\hyperref[task:4]{4} &
Прямая $y=kx$ и выколотая точка &
При сокращении потерян минус перед единицей: вместо $-x^2-1$ получено $-x^2+1$. Из-за этого дальнейший график и значение $k$ неверны; кроме того, не рассмотрен случай касания. &
$k=-\frac{5}{3}$. &
$k=-2$, $k=2$, $k=-\frac{10}{3}$. \\
\midrule
\hyperref[task:5]{5} &
Прямая $y=kx$ и выколотая точка &
Решение не выполнено. &
Ответ не указан. &
$k=-6$, $k=6$, $k=\frac{13}{2}$. \\
\bottomrule
\end{longtable}
\normalsize

\newpage
\thispagestyle{empty}
\vspace*{0.38\textheight}
\begin{center}
{\Huge\bfseries Разбор ошибок}
\end{center}
\vfill

% ------------------------------------------------------------
\newpage
\phantomsection\label{task:2}
\section*{Задание 2}

\subsection*{Текст задания}
Построить график функции
\[
y=\frac{(x^2+4)(x-1)}{1-x}
\]
и указать вершину параболы, пересечения с осями и выколотую точку.

\subsection*{Ваше решение}
В решении верно замечено, что
\[
1-x=-(x-1),
\]
поэтому после сокращения получено
\[
y=-(x^2+4)=-x^2-4.
\]
Далее найдена вершина с абсциссой $0$ и ординатой $-4$, установлено отсутствие пересечений с осью $Ox$, а пересечение с осью $Oy$ получено в точке $(0,-4)$. Затем составлена таблица значений, в которой присутствует значение $x=1$, и на координатной плоскости нанесены отдельные точки.

\subsection*{Ваш ответ}
Вершина $(0,-4)$, пересечений с осью $Ox$ нет, пересечение с осью $Oy$ --- $(0,-4)$. Координаты выколотой точки отдельно не указаны; полноценный график параболы не построен.

\subsection*{Ошибка}
\errlabel\ значение $x=1$ использовано как обычная точка графика, хотя оно исключено из области определения исходной функции.

\subsection*{Где именно возникает ошибка}
Последний полностью правильный шаг:
\[
y=-x^2-4.
\]
Но это равенство описывает исходную функцию только при условии
\[
x\ne1.
\]
Первый неверный или неполный шаг появляется при составлении таблицы: для $x=1$ записывается значение $y=-5$ как для обычной точки графика. Именно эту точку исходная функция не содержит.

\subsection*{Почему это неверно}
В исходной дроби при $x=1$ знаменатель равен нулю:
\[
1-x=0.
\]
Следовательно, исходное выражение при $x=1$ не определено. Сокращение множителя $(x-1)$ упрощает формулу, но не возвращает исключённое значение в область определения. Поэтому функция совпадает с параболой
\[
y=-x^2-4
\]
во всех точках, кроме точки с абсциссой $1$.

Если формально подставить $x=1$ в упрощённое выражение, получится
\[
y=-1-4=-5.
\]
Именно поэтому на параболе нужно удалить точку $(1,-5)$ и обозначить её пустым кружком.

\subsection*{Ключевая идея}
\keyidea\ после сокращения дроби обязательно сохраняется ограничение, пришедшее из исходного знаменателя. Упрощённая формула задаёт форму графика, а исходная область определения показывает, какие точки из этого графика нужно удалить.

\subsection*{Исправление от последнего правильного шага}
Оставляем верное упрощение
\[
y=-x^2-4,
\]
но добавляем ограничение
\[
x\ne1.
\]
После этого значение $x=1$ не используем как обычную точку. Вместо этого вычисляем координаты удаляемой точки по упрощённой формуле:
\[
y(1)=-1^2-4=-5.
\]
Значит, точка $(1,-5)$ должна быть выколота.

\subsection*{Полное правильное решение}
Сначала определим область определения исходной функции. Знаменатель не должен обращаться в ноль:
\[
1-x\ne0,
\]
откуда
\[
x\ne1.
\]

Преобразуем функцию:
\[
\frac{(x^2+4)(x-1)}{1-x}
=
\frac{(x^2+4)(x-1)}{-(x-1)}
=-(x^2+4)
=-x^2-4,
\]
но только при $x\ne1$.

Получена парабола с ветвями, направленными вниз. Для функции
\[
y=-x^2-4
\]
коэффициент при $x^2$ отрицателен, а линейного члена нет, поэтому вершина имеет абсциссу
\[
x_v=0.
\]
Тогда
\[
y_v=-0^2-4=-4.
\]
Вершина:
\[
(0,-4).
\]

Пересечения с осью $Ox$ ищем из условия $y=0$:
\[
-x^2-4=0,
\]
\[
x^2=-4.
\]
Действительных решений нет, поэтому с осью $Ox$ график не пересекается.

Пересечение с осью $Oy$ получаем при $x=0$:
\[
y=-4.
\]
Точка пересечения:
\[
(0,-4).
\]

Выколотая точка соответствует исключённому значению $x=1$. По упрощённой формуле её ордината равна
\[
-1^2-4=-5.
\]
Следовательно, выколотая точка:
\[
(1,-5).
\]

Для построения нужно провести параболу $y=-x^2-4$ с вершиной $(0,-4)$ и вместо обычной точки $(1,-5)$ поставить пустой кружок.

\subsection*{Проверка}
Проверим ключевые характеристики. При $x=0$ получаем $y=-4$, поэтому вершина и пересечение с осью $Oy$ согласованы. Значение $-x^2-4$ всегда отрицательно, поэтому пересечений с осью $Ox$ действительно нет. При $x=1$ исходный знаменатель равен нулю, поэтому точка $(1,-5)$ действительно должна быть исключена.

\subsection*{Итог}
\resultlabel\ вычислительная часть почти полностью выполнена верно. Исправить нужно именно учёт области определения при построении: график --- парабола $y=-x^2-4$ с выколотой точкой $(1,-5)$.

\subsection*{Задача на закрепление}
Постройте график функции
\[
y=\frac{(x^2+1)(x-2)}{2-x}.
\]
Укажите вершину параболы, пересечения с осями и выколотую точку.

% ------------------------------------------------------------
\newpage
\phantomsection\label{task:3}
\section*{Задание 3}

\subsection*{Текст задания}
Для функции из задания 1 определить количество общих точек её графика с прямой
\[
y=m
\]
в зависимости от значения $m$.

\subsection*{Ваше решение}
Записано, что одна общая точка получается при $m=-6{,}25$ и при $m=-4$, две общие точки --- при $m>-4$, а общих точек нет при $m<-6{,}25$.

\subsection*{Ваш ответ}
Одна точка при $m=-\dfrac{25}{4}$ и $m=-4$; две точки при $m>-4$; нет общих точек при $m<-\dfrac{25}{4}$.

\subsection*{Ошибка}
\errlabel\ пропущен промежуток
\[
-\frac{25}{4}<m<-4.
\]
На всём этом промежутке горизонтальная прямая также пересекает график в двух точках.

\subsection*{Где именно возникает ошибка}
Верно определены два особых уровня:
\[
m=-\frac{25}{4}
\]
--- уровень вершины параболы, и
\[
m=-4
\]
--- уровень выколотой точки.

Первый неполный шаг --- утверждение, что две общие точки существуют только при $m>-4$. Между уровнем вершины и уровнем выколотой точки, то есть при
\[
-\frac{25}{4}<m<-4,
\]
прямая $y=m$ также пересекает обе ветви параболы, и ни одна из этих точек не совпадает с выколотой точкой.

\subsection*{Почему это неверно}
Из задания 1 после сокращения получается
\[
y=x^2-x-6,
\]
причём
\[
x\ne2.
\]
Это парабола с ветвями вверх. Её минимальное значение равно
\[
-\frac{25}{4}.
\]
Поэтому для любой горизонтальной прямой, расположенной строго выше вершины, у полной параболы имеются две точки пересечения.

Исключение возникает только на уровне выколотой точки. При $x=2$ упрощённая формула даёт
\[
y=2^2-2-6=-4.
\]
Значит, на уровне $m=-4$ одна из двух точек полной параболы удалена, и остаётся ровно одна общая точка.

Во всех остальных уровнях выше вершины обе точки сохраняются. Поэтому промежуток
\[
-\frac{25}{4}<m<-4
\]
нельзя пропускать.

\subsection*{Ключевая идея}
\keyidea\ сначала учитывается обычное число пересечений горизонтальной прямой с параболой, а затем отдельно проверяется, не совпадает ли одна из найденных точек с выколотой точкой.

\subsection*{Исправление от последнего правильного шага}
Сохраняем верные особые значения $m=-\dfrac{25}{4}$ и $m=-4$, но разбиваем все значения $m$ полностью:
\begin{itemize}
\item ниже вершины --- пересечений нет;
\item на уровне вершины --- одна точка;
\item выше вершины --- обычно две точки;
\item на уровне выколотой точки $m=-4$ одна из двух точек удаляется, поэтому остаётся одна.
\end{itemize}

\subsection*{Полное правильное решение}
Функция из задания 1 после сокращения имеет вид
\[
y=x^2-x-6,
\]
при условии
\[
x\ne2.
\]

Ищем общие точки с прямой $y=m$:
\[
x^2-x-6=m.
\]
Перенесём всё в левую часть:
\[
x^2-x-(m+6)=0.
\]
Дискриминант этого квадратного уравнения равен
\[
D=(-1)^2-4\cdot1\cdot(-(m+6))=1+4m+24=4m+25.
\]

Если
\[
4m+25<0,
\]
то
\[
m<-\frac{25}{4},
\]
и действительных корней нет. Значит, общих точек нет.

Если
\[
4m+25=0,
\]
то
\[
m=-\frac{25}{4}.
\]
Уравнение имеет один двойной корень, соответствующий вершине параболы. Эта точка не выколота, поэтому общая точка ровно одна.

Если
\[
4m+25>0,
\]
то
\[
m>-\frac{25}{4},
\]
и у полной параболы две точки пересечения. Теперь проверяем выколотую точку. Исключённое значение $x=2$ соответствует уровню
\[
m=2^2-2-6=-4.
\]
При $m=-4$ уравнение принимает вид
\[
x^2-x-2=0,
\]
то есть
\[
(x-2)(x+1)=0.
\]
Корни равны $2$ и $-1$, но $x=2$ исключён. Поэтому при $m=-4$ остаётся только одна общая точка.

При любом другом значении
\[
m>-\frac{25}{4},\qquad m\ne-4,
\]
оба корня допустимы, и общих точек две.

Итак,
\[
\begin{cases}
0\text{ общих точек}, & m<-\dfrac{25}{4},\\[4pt]
1\text{ общая точка}, & m=-\dfrac{25}{4}\text{ или }m=-4,\\[4pt]
2\text{ общие точки}, & m>-\dfrac{25}{4},\ m\ne-4.
\end{cases}
\]

\subsection*{Проверка}
Возьмём значение из пропущенного промежутка, например $m=-5$. Тогда
\[
x^2-x-1=0.
\]
Дискриминант равен $5>0$, поэтому имеются два действительных корня. Ни один из них не равен $2$, следовательно, обе точки принадлежат исходному графику. Это подтверждает, что при $-\dfrac{25}{4}<m<-4$ общих точек две.

\subsection*{Итог}
\resultlabel\ правильная классификация должна учитывать весь диапазон значений выше минимума параболы. Единственное дополнительное исключение --- уровень $m=-4$, где одна из двух точек является выколотой.

\subsection*{Задача на закрепление}
Для функции
\[
y=\frac{(x+1)(x^2-4x+3)}{x-3}
\]
определите количество общих точек её графика с прямой $y=m$ в зависимости от значения $m$.

% ------------------------------------------------------------
\newpage
\phantomsection\label{task:4}
\section*{Задание 4}

\subsection*{Текст задания}
Определить все значения $k$, при которых прямая
\[
y=kx
\]
имеет ровно одну общую точку с графиком функции
\[
y=\frac{(x^2+1)(x-3)}{3-x}.
\]

\subsection*{Ваше решение}
Начато сокращение дроби, после которого записано
\[
y=-x^2+1.
\]
Далее построена парабола с вершиной выше оси $Ox$, отмечены точки пересечения с осью $Ox$, затем по выбранной точке получено равенство $-5=3k$ и ответ
\[
k=-\frac{5}{3}.
\]

\subsection*{Ваш ответ}
\[
k=-\frac{5}{3}.
\]

\subsection*{Ошибка}
\errlabel\ при раскрытии минуса потерян знак перед единицей:
\[
-(x^2+1)=-x^2-1,
\]
а не $-x^2+1$.

\subsection*{Где именно возникает ошибка}
Последний правильный шаг состоит в том, что
\[
3-x=-(x-3).
\]
Поэтому при $x\ne3$
\[
\frac{(x^2+1)(x-3)}{3-x}
=
\frac{(x^2+1)(x-3)}{-(x-3)}
=-(x^2+1).
\]

Первый неверный переход:
\[
-(x^2+1)\longrightarrow -x^2+1.
\]
Знак минус перед скобками должен изменить знак каждого слагаемого внутри скобок. Поэтому верно
\[
-(x^2+1)=-x^2-1.
\]

\subsection*{Почему это неверно}
Правило раскрытия скобок с минусом означает умножение каждого слагаемого на $-1$:
\[
-1\cdot x^2=-x^2,
\qquad
-1\cdot1=-1.
\]
После ошибки функция была заменена другой параболой. Из-за этого изменились вершина, пересечения с осями, положение выколотой точки и все дальнейшие рассуждения о прямой $y=kx$.

Кроме того, даже после правильного построения нельзя искать только прямую, проходящую через выколотую точку. Ровно одна общая точка возникает ещё и тогда, когда прямая $y=kx$ касается параболы, то есть соответствующее квадратное уравнение имеет один двойной корень.

\subsection*{Ключевая идея}
\keyidea\ условие «ровно одна общая точка» здесь возникает двумя разными способами: либо прямая касается параболы, либо у полной параболы имеются две точки пересечения, но одна из них совпадает с выколотой точкой и не принадлежит исходному графику.

\subsection*{Исправление от последнего правильного шага}
От правильного сокращения получаем
\[
y=-x^2-1,
\qquad x\ne3.
\]
Теперь вместо чтения коэффициента $k$ с рисунка составляем уравнение пересечения с прямой $y=kx$:
\[
-x^2-1=kx.
\]
После переноса:
\[
x^2+kx+1=0.
\]
Далее отдельно рассматриваем случай двойного корня и случай, когда одним из двух корней является запрещённое значение $x=3$.

\subsection*{Полное правильное решение}
Сначала упростим функцию. Область определения задаётся условием
\[
3-x\ne0,
\]
то есть
\[
x\ne3.
\]
Так как
\[
3-x=-(x-3),
\]
то при $x\ne3$
\[
y=\frac{(x^2+1)(x-3)}{3-x}
=-(x^2+1)
=-x^2-1.
\]
В точке $x=3$ исходная функция не определена. По упрощённой формуле соответствующая ордината равна
\[
-3^2-1=-10,
\]
поэтому выколотая точка имеет координаты
\[
(3,-10).
\]

Теперь ищем пересечения с прямой $y=kx$:
\[
-x^2-1=kx.
\]
Переносим все слагаемые в одну сторону:
\[
x^2+kx+1=0.
\]

\textbf{Первый способ получить ровно одну общую точку: касание.}
Квадратное уравнение должно иметь один двойной корень, то есть
\[
D=k^2-4=0.
\]
Отсюда
\[
k^2=4,
\]
\[
k=2\quad\text{или}\quad k=-2.
\]
При $k=2$ двойной корень равен $x=-1$, а при $k=-2$ --- $x=1$. Оба значения допустимы, так как не равны $3$. Следовательно, оба коэффициента подходят.

\textbf{Второй способ: одна из двух точек пересечения выколота.}
Для этого запрещённое значение $x=3$ должно быть корнем уравнения
\[
x^2+kx+1=0.
\]
Подставим $x=3$:
\[
9+3k+1=0,
\]
\[
3k+10=0,
\]
\[
k=-\frac{10}{3}.
\]
При этом значение дискриминанта положительно:
\[
D=\frac{100}{9}-4=\frac{64}{9}>0,
\]
поэтому у полной параболы действительно две точки пересечения. Один корень равен $3$ и удалён из графика. По теореме Виета произведение корней равно $1$, значит второй корень равен
\[
\frac{1}{3},
\]
и он допустим. Следовательно, при $k=-\dfrac{10}{3}$ остаётся ровно одна общая точка.

Итак, все подходящие значения:
\[
\boxed{k=-2,\quad k=2,\quad k=-\frac{10}{3}}.
\]

\subsection*{Проверка}
Для $k=2$ уравнение имеет вид
\[
x^2+2x+1=(x+1)^2=0,
\]
то есть одна точка с $x=-1$.

Для $k=-2$:
\[
x^2-2x+1=(x-1)^2=0,
\]
то есть одна точка с $x=1$.

Для $k=-\dfrac{10}{3}$ корни равны $3$ и $\dfrac13$. Первый корень соответствует выколотой точке и не учитывается, второй остаётся. Значит, во всех трёх случаях общая точка действительно ровно одна.

\subsection*{Итог}
\resultlabel\ исходная ошибка --- неверное раскрытие минуса. После исправления нужно учитывать два механизма получения единственной общей точки: касание и удаление одной из двух точек пересечения из-за ограничения $x\ne3$.

\subsection*{Задача на закрепление}
Определите все значения $k$, при которых прямая $y=kx$ имеет ровно одну общую точку с графиком функции
\[
y=\frac{(x^2+1)(x-2)}{2-x}.
\]

% ------------------------------------------------------------
\newpage
\phantomsection\label{task:5}
\section*{Задание 5}

\subsection*{Текст задания}
Определить все значения $k$, при которых прямая
\[
y=kx
\]
имеет ровно одну общую точку с графиком функции
\[
y=\frac{(x^2+9)(x+2)}{-x-2}.
\]

\subsection*{Ваше решение}
Решение не выполнено. В работе записано только обозначение задания и знак вопроса.

\subsection*{Ваш ответ}
Не указан.

\subsection*{Ошибка}
\errlabel\ задание оставлено без решения, поэтому не выполнены ни упрощение функции, ни учёт выколотой точки, ни анализ числа пересечений с прямой $y=kx$.

\subsection*{Где именно возникает ошибка}
Последнего выполненного математического шага нет. Решение должно начинаться с области определения и сокращения общего множителя $x+2$ с обязательным сохранением ограничения $x\ne-2$.

\subsection*{Почему это важно}
В этой задаче недостаточно рассмотреть только касание прямой и параболы. После сокращения в графике остаётся выколотая точка. Поэтому ровно одна общая точка может возникнуть и в ситуации, когда у полной параболы две точки пересечения с прямой, но одна из них является исключённой точкой исходной функции.

\subsection*{Ключевая идея}
\keyidea\ сначала получить простую квадратичную функцию с сохранением ограничения из знаменателя, затем решить уравнение пересечения с $y=kx$ и отдельно проверить два случая: двойной корень и запрещённый корень.

\subsection*{Исправление}
Начинаем решение с преобразования знаменателя:
\[
-x-2=-(x+2).
\]
При $x\ne-2$ общий множитель $x+2$ сокращается, но значение $x=-2$ остаётся запрещённым.

\subsection*{Полное правильное решение}
Определим область определения:
\[
-x-2\ne0,
\]
откуда
\[
x\ne-2.
\]

Упростим функцию:
\[
y=\frac{(x^2+9)(x+2)}{-(x+2)}
=-(x^2+9)
=-x^2-9,
\]
при условии
\[
x\ne-2.
\]

Найдём выколотую точку. По упрощённой формуле при $x=-2$ получаем
\[
y=-(-2)^2-9=-4-9=-13.
\]
Значит, выколотая точка:
\[
(-2,-13).
\]

Теперь ищем пересечения графика с прямой $y=kx$:
\[
-x^2-9=kx.
\]
Переносим все слагаемые:
\[
x^2+kx+9=0.
\]

\textbf{Первый способ получить ровно одну общую точку: касание.}
Для одного двойного корня нужно
\[
D=k^2-36=0.
\]
Следовательно,
\[
k^2=36,
\]
\[
k=6\quad\text{или}\quad k=-6.
\]
При $k=6$ двойной корень равен $x=-3$, а при $k=-6$ --- $x=3$. Оба значения допустимы, так как не равны $-2$. Значит, оба коэффициента подходят.

\textbf{Второй способ: один из двух корней запрещён.}
Требуем, чтобы $x=-2$ был корнем уравнения
\[
x^2+kx+9=0.
\]
Подставляем:
\[
4-2k+9=0,
\]
\[
13-2k=0,
\]
\[
k=\frac{13}{2}.
\]
Для этого значения дискриминант положителен:
\[
D=\frac{169}{4}-36=\frac{25}{4}>0,
\]
значит, у полной параболы две точки пересечения. Один корень равен $-2$ и не принадлежит исходному графику. По теореме Виета произведение корней равно $9$, поэтому второй корень равен
\[
\frac{9}{-2}=-\frac{9}{2},
\]
и он допустим. Значит, после удаления выколотой точки остаётся ровно одна общая точка.

Итак,
\[
\boxed{k=-6,\quad k=6,\quad k=\frac{13}{2}}.
\]

\subsection*{Проверка}
При $k=6$:
\[
x^2+6x+9=(x+3)^2=0,
\]
поэтому имеется одна допустимая точка.

При $k=-6$:
\[
x^2-6x+9=(x-3)^2=0,
\]
поэтому также имеется одна допустимая точка.

При $k=\dfrac{13}{2}$ корни равны $-2$ и $-\dfrac92$. Корень $-2$ исключён исходной областью определения, а второй корень допустим. Поэтому и в этом случае общая точка ровно одна.

\subsection*{Итог}
\resultlabel\ задание имеет три ответа. Два из них возникают из условия касания, а третий --- из-за того, что одна из двух точек пересечения совпадает с выколотой точкой исходного графика.

\subsection*{Задача на закрепление}
Определите все значения $k$, при которых прямая $y=kx$ имеет ровно одну общую точку с графиком функции
\[
y=\frac{(x^2+4)(x+3)}{-x-3}.
\]

\vfill
\begin{center}
\small Лёвин Артём Александрович эксклюзивно для Марины Селиверстовой
\end{center}

\end{document}
